\subsection*{引理 5.4.}
\underline{假设 $X$ 的某个光滑超平面截面 (因而由推论 5.2.7, 每个这样的截面) 满足 (LV)---例如, 假设 $X$ 满足 (LV), 参见引理 5.2.4---并假设 $\rho$ 满足条件 (A). 则每个纤维 (包括每个奇异纤维) 都满足 (LV), 且当 $i\neq n-1$ 时, 层 $R^i\rho_*\mathbb Q_\ell(j)$ 为常层.}

\underline{证明.}
设
\[
L:R^i\rho_*\mathbb Q_\ell(j)
\longrightarrow R^{i+2}\rho_*\mathbb Q_\ell(j+1)
\]
表示与 $f^*c_1(\mathcal O_X(1))\in H^2(\widetilde X,\mathbb Z_\ell(1))$ 相乘的映射; 逐纤维来看, 这就是 (5.2.1) 中考虑的运算. 由假设, 在 $U$ 上有同构
\begin{equation}\tag{5.4.1}
\nu^*(L^i):
\nu^*R^{n-1-i}\rho_*\mathbb Q_\ell(j)
\overset{\sim}{\longrightarrow}
\nu^*R^{n-1+i}\rho_*\mathbb Q_\ell(j+i)
\end{equation}
对每个 $i>0$ 和每个 $j$ 都成立. 应用函子 $\nu_*$, 得到\underline{同构}
\begin{equation}\tag{5.4.2}
(\nu_*\nu^*)L^i:
\nu_*\nu^*R^{n-1-i}\rho_*\mathbb Q_\ell(j)
\overset{\sim}{\longrightarrow}
\nu_*\nu^*R^{n-1+i}\rho_*\mathbb Q_\ell(j+i).
\end{equation}
由 (A), 它们可与下列态射对应:
\begin{equation}\tag{5.4.3}
L^i:R^{n-1-i}\rho_*\mathbb Q_\ell(j)
\overset{\sim}{\longrightarrow}
R^{n-1+i}\rho_*\mathbb Q_\ell(j+i).
\end{equation}
结论由 5.3.3 得出.

\subsubsection*{5.4.4.}
设 $Y$ 为 $\rho:\widetilde X\to\mathbb P^1$ 的一个光滑纤维. 由弱 Lefschetz 定理, 限制映射
\[
H^{n-1}(X,\mathbb Q_\ell(j))
\overset{m^*}{\hookrightarrow}
H^{n-1}(Y,\mathbb Q_\ell(j))
\]
为单射. 我们已在 5.2.4 中看到, 若 $X$ 满足 (LV), 则 $H^{n-1}(Y,\mathbb Q_\ell(j))$ 上的杯积限制到 $m^*H^{n-1}(X,\mathbb Q_\ell(j))$ 后仍非退化. 因而 $H^{n-1}(X,\mathbb Q_\ell(j))$ 在 $H^{n-1}(Y,\mathbb Q_\ell(j))$ 中的正交补, 是 $H^{n-1}(X,\mathbb Q_\ell(j))$ 在 $H^{n-1}(Y,\mathbb Q_\ell(j))$ 中的一个补子空间; 依循经典术语, 称之为 \underline{$Y$ 的消灭上同调空间}.

设 $\nu:U\hookrightarrow\mathbb P^1$ 为任意非空开子集, 使 $\rho$ 在其上光滑. $\nu_*\nu^*R^{n-1}\rho_*\mathbb Q_\ell(j)$ 的相应子层, 即 $\mathbb P^1$ 上的“\underline{消灭上同调层}” $E^{n-1}\rho_*\mathbb Q_\ell(j)$, 定义为
\begin{equation}\tag{5.4.4.1}
E^{n-1}\rho_*\mathbb Q_\ell(j)
=\nu_*\left\{
\begin{array}{l}
\text{在 }
\nu^*R^{n-1}\rho_*\mathbb Q_\ell(j)\text{ 中}\\[1ex]
\text{下列常子层的正交补: }H^{n-1}(X,\mathbb Q_\ell(j))_U
\end{array}
\right\}.
\end{equation}
显然, 此层不依赖于 $U$ 的选取.

\subsection*{引理 5.5.}
\underline{假设 $X$ 满足 (LV), 且当 $i=n-1$ 时映射 (5.3.2) 为同构---例如, 假设 $\rho$ 来自 Lefschetz 线束, 参见 5.3.4. 则有直和分解}
\begin{equation}\tag{5.5.1}
R^{n-1}\rho_*\mathbb Q_\ell(j)
\simeq
H^{n-1}(X,\mathbb Q_\ell(j))_{\mathbb P^1}
\oplus E^{n-1}\rho_*\mathbb Q_\ell(j).
\end{equation}

\underline{证明.}
由 (LV), 在 $U$ 上有分解
\begin{equation}\tag{5.5.2}
\nu^*R^{n-1}\rho_*\mathbb Q_\ell(j)
\simeq
H^{n-1}(X,\mathbb Q_\ell(j))_U
\oplus \nu^*E^{n-1}\rho_*\mathbb Q_\ell(j),
\end{equation}
对它应用 $\nu_*$, 即得所断言的结论.
